\section{Purpose, Scope, and Properties}
\label{sec:sec01}
\label{sec:intro}

\subsection{The Problem}

Conventional post-trade infrastructure records the same economic reality in many places: trading systems, risk engines, general ledgers, custodian records, regulatory reports. Each carries its own copy of positions, valuations, and histories. Reconciliation---verifying that these independent records agree---is the dominant operational cost and the primary source of breaks. The cause is architectural: a multi-source-of-truth design admits inconsistency by construction.

\subsection{The Solution}

The post-trade activities within scope use a single primitive: the atomic move of a quantity between wallets. Financial instruments---equities, bonds, derivatives, swaps, managed accounts, QIS strategies---are deterministic contracts that emit sequences of moves. An immutable move stream records every move with full provenance and is the sole canonical record; every other view---balances, profit and loss, balance sheets, reports---is a projection of it.

The framework satisfies six properties by construction.

\begin{enumerate}
\item \textbf{Atomicity.} Every transaction is an indivisible operation with complete provenance.

\item \textbf{Conservation.} The system is closed: the \emph{quantity} of each unit is neither created nor destroyed. This is double-entry accounting reduced to one equation, $\sum_w w(u) = 0$. Conservation governs quantities per unit, not market values; value changes from price movements are captured by the valuation layer (Section~\ref{sec:sec05}).

\item \textbf{Determinism.} Given inputs and conditions, execution is reproducible.

\item \textbf{State-Sufficiency.} Portfolio value depends only on current state---wallet balances, unit state (including embedded lifecycle state: accrued coupons, exercise status, barrier flags), and the current price vector---not on transaction history. Where unit state lives is fixed by the three-home model (Section~\ref{sec:sec04}).

\item \textbf{Lifecycle Value Invariance} (restricted form). \label{prop:lifecycle-value-invariance} {\setlength{\emergencystretch}{2em}For contractual lifecycle events involving only deterministic cash flows---\allowbreak coupon and dividend payments, resets, principal redemptions---\allowbreak the total value of all wallets under a given price vector is invariant: the jump in the instrument's price is exactly offset by the explicit cash move. Lifecycle processing reshuffles value between units and cash; it creates and destroys no PnL. This connects conservation, which governs quantities, to value: across deterministic lifecycle events total value is also invariant.\par}

\textit{Qualification.} For lifecycle events involving optionality---exercise, expiry, barrier knock-in/out, callable redemption---the invariance holds only for intrinsic value. Time value and embedded option value are extinguished when optionality resolves: exercising an option destroys time value; calling a bond extinguishes the holder's above-market coupon stream. This is not a conservation violation, because those components reflect probabilistic expectations that cease to exist when uncertainty resolves. Conservation governs quantities, not market values; value invariance holds exactly for deterministic events and for intrinsic value at option-exercise events.

\item \textbf{Time Travel.} The exact ledger state at any historical timestamp is reconstructed by replaying the immutable move stream and unit-state events up to that time. The operation \texttt{clone\_at($t$)} yields the ledger as of time $t$, replaying activity up to $t$ and ignoring the rest. Combined with pure lifecycle functions, this yields historical truth: the ledger is restored to what was known at time $t$ and re-simulated. ``Time travel to what was known at $t$'' (market data as captured at $t$) and ``time travel to $t$ with today's corrected data'' (restated data) are distinct and are not conflated.
\end{enumerate}

These properties make the ledger a deterministic function of its event history: given the same transactions and market data, the state at any time is uniquely determined.

\subsection{Three Unreachable Internal Breaks}

When the ledger is the sole system of record for the activities it covers, three of the most damaging reconciliation failures are structurally impossible within that scope:

\begin{enumerate}[nosep]
\item trading-system / general-ledger divergence,
\item PnL / balance-sheet mismatch,
\item inter-entity breaks.
\end{enumerate}

\noindent Each is a disagreement between views that are, by construction, projections of one record; disagreement is not representable. The guarantee is scoped: during any migration in which legacy systems run in parallel, it covers only activity recorded in the ledger. Activity remaining in legacy systems retains the legacy failure modes.

\subsection{Architecture}

A \textit{wallet} holds quantities of \textit{units} (currencies, securities, derivatives, or any contractual obligation). A \textit{move} transfers a quantity of a unit between two wallets. A \textit{transaction} is a finite set of simultaneous moves satisfying conservation. A \textit{smart contract} is a deterministic program mapping inputs and state to sequences of moves. Unit state records each unit's lifecycle position. The immutable \textit{move stream} records every move with full provenance; all other views are derived projections. External authorities---CSDs, share registrars, regulators---maintain independent records at the boundary.

\subsection{Scope}

The scope is trading books measured at fair value through profit or loss under IFRS and US~GAAP.\footnote{IFRS~9 \S\S4.1.2, 4.1.4, 5.7.1--5.7.7; ASC~815, ASC~320/321, ASC~820.} Positions are marked to market using price vectors consistent with IFRS~13 and ASC~820 fair-value guidance. Banking-book activities---amortised-cost loans, held-to-collect debt, hedge accounting for structural positions---are out of scope and extend the same primitives.

The ISDA Common Domain Model (CDM) is the canonical vocabulary: it defines the unit universe $\mathcal{U}$, drives lifecycle state transitions, bounds the input space for property-based testing, and maps ledger transactions to ISO~20022 settlement messages. CDM integration, the product hierarchy, and its layer-by-layer justification are developed in Section~\ref{sec:sec13} and Appendix~\ref{app:cdm-walkthrough}.

\subsection{The Ledger Boundary}
\label{sec:scope-overview}

The framework's guarantees hold for activity recorded inside the ledger boundary. The boundary is stated here, once.

\medskip
\noindent\textbf{Inside} (governed by the framework's invariants):
\begin{itemize}[nosep]
\item Positions: wallet balances, the generalised position vector, the conservation law $\sum_w w_t(u) = 0$.
\item Moves: atomic transfers of quantity between wallets, with full provenance.
\item Lifecycle events: state transitions on units, recorded as transactions with CDM event metadata.
\item PnL: portfolio valuation $V_t = \sum_w w_t(u)\,P_t(u)$ for an externally supplied price vector $P_t$.
\item Time travel: reconstruction of historical state by move-stream and unit-state replay.
\item Internal reconciliation: the guarantee that trading-system, general-ledger, and balance-sheet views are projections of one record.
\end{itemize}

\medskip
\noindent\textbf{Outside} (external systems, not governed by the framework):
\begin{itemize}[nosep]
\item Price formation: $P_t$ is an external input; the framework does not compute, validate, or guarantee executable prices.
\item Legal agreements: ISDA Master Agreements, GMSLAs, GMRAs, CSAs, and custody agreements determine legal rights. The ledger records economic positions, not legal ownership or enforceability.
\item Settlement infrastructure: CSD connectivity, payment-system access, settlement finality, and real-world delivery-versus-payment (DvP) depend on external infrastructure (DTC, Euroclear, Clearstream, CLS, TARGET2).
\item Regulatory gateways: trade, transaction, and prudential reporting (EMIR, MiFIR, SFTR, Dodd-Frank, MiFID~II RTS~25, CRR, FRTB) require submission infrastructure beyond the ledger.
\item Reference-data authority: instrument master data (ISINs, contract specifications, corporate-action terms) is consumed from external sources, not created by the ledger.
\item Model governance: pricing and risk models and their validation.
\item Accounting policy: IFRS~9 classification (FVTPL, FVOCI, amortised cost), hedge-accounting designation, and impairment assessment.
\end{itemize}

\noindent Every ``within scope'' qualifier elsewhere in this document refers to the inside list above.
